\icode{E-T-CA-L3-261011-001}
如图所示，一导热性能良好的球形容器内部不规则，某兴趣小组为了测量它的容积，在容器上竖直插入一根两端开口的长玻璃管，接口密封，玻璃管内部横截面积 \( S = \) \,\qty{0.2}{cm^2}，一长为 \( h = \) \,\qty{15}{cm}的静止水银柱封闭了一定质量的气体，其下方玻璃管内空气柱的长度 \( l_1 = \) \,\qty{10}{cm}，此时外界温度 \( t_1 = \) \,\qty{27}{\degreeCelsius}。现把容器浸在\qty{77}{\degreeCelsius}的热水中，水银柱缓慢上升\qty{30}{cm}后稳定。实验过程中认为大气压没有变化，大气压 \( p_0 = \) \,\qty{1.0e5}{Pa}（相当于\qty{75}{cm}高汞柱压强）。
\begin{enumerate}
    \item 容器的容积为多少？
    \item 若实验过程中管内气体的内能增加了\qty{1.5}{J}。判断气体是从外界吸收热量还是向外界放出热量，并计算热量的大小。
\end{enumerate}

\begin{tikzpicture}
    \draw (97:.5) arc(97:443:.5);
    \draw[fill=gray] ($(97:.5)+(0,.8)$) rectangle ($(83:.5)+(0,1.5)$);
    \draw (97:.5)--++(0,2.4) (83:.5)--++(0,2.4);
\end{tikzpicture}